\documentclass[11pt]{article}
\usepackage[margin=1in]{geometry}
\usepackage{xcolor}
\usepackage[breakable,listings]{tcolorbox}
\tcbuselibrary{listings}
\usepackage{hyperref}
\usepackage{enumitem}
\usepackage{listings}
\usepackage{amsmath}
\usepackage{booktabs}
\usepackage{array}

% Define colors
\definecolor{osaorange}{RGB}{220,115,38}
\definecolor{aiblue}{RGB}{30,144,255}
\definecolor{answergreen}{RGB}{34,139,34}

% Configure hyperref
\hypersetup{
    colorlinks=true,
    linkcolor=blue,
    urlcolor=blue
}

% Code environment
\newtcblisting{rcode}{
  colback=white,
  colframe=gray!50,
  boxrule=0.5pt,
  arc=0pt,
  left=3pt,
  right=3pt,
  top=3pt,
  bottom=3pt,
  width=\textwidth,
  breakable,
  listing only,
  listing options={
    basicstyle=\small\ttfamily,
    breaklines=true,
    breakatwhitespace=true,
    columns=flexible,
    keepspaces=true,
    showstringspaces=false
  }
}

\title{}
\author{}
\date{}

\begin{document}

% Header box
\begin{tcolorbox}[colback=osaorange,colframe=black,boxrule=2pt,arc=0pt,left=3pt,right=3pt,top=6pt,bottom=6pt,boxsep=2pt]
\centering
\textcolor{white}{\textbf{\Large H524 Introduction to Biostatistics}}\\
\textcolor{white}{\textbf{\Large Week 6: Power, Two-Sample Tests, and ANOVA}}\\
\textcolor{white}{\textbf{\Large Homework Assignment - ANSWER KEY}}\\
\textcolor{white}{\textbf{\Large Fall 2025}}
\end{tcolorbox}

\vspace{8pt}

\noindent\textbf{Note:} All numerical answers verified using R (see \texttt{verify\_assignment\_answers.R})

\vspace{12pt}

\section{Numerical Questions (10 points total)}

\subsection*{Question 1: Sample Size for Power (5 points)}

\textbf{Problem:} Calculate the required sample size to detect a 10 mmHg reduction in blood pressure with 80\% power, using $\sigma = 20$ mmHg and $\alpha = 0.05$ (two-sided).

\vspace{0.3cm}

\textbf{R Code:}

\begin{rcode}
# Sample size calculation for one-sample t-test
result <- power.t.test(delta = 10,       # Effect size
                       sd = 20,           # Standard deviation
                       sig.level = 0.05,  # Alpha
                       power = 0.80,      # Desired power
                       type = "one.sample",
                       alternative = "two.sided")
print(result)

# Round up to nearest whole number
n_required <- ceiling(result$n)
cat("Required sample size:", n_required, "subjects\n")
\end{rcode}

\vspace{0.3cm}

\textbf{Output:}
\begin{verbatim}
     One-sample t test power calculation

              n = 33.37
          delta = 10
             sd = 20
      sig.level = 0.05
          power = 0.8
    alternative = two.sided

Required sample size: 34 subjects
\end{verbatim}

\vspace{0.3cm}

\textcolor{answergreen}{\textbf{ANSWER: 34 subjects}}

\vspace{0.3cm}

\textbf{Interpretation:} To have an 80\% chance of detecting a true reduction of 10 mmHg in blood pressure (with $\alpha = 0.05$), the study needs to enroll at least 34 subjects.

\vspace{1cm}

\subsection*{Question 2: Two-Sample t-Test and ANOVA (5 points)}

\textbf{Data:}
\begin{center}
\begin{tabular}{lll}
\toprule
Program A & Program B & Program C \\
\midrule
4.2 & 6.1 & 7.8 \\
3.8 & 5.8 & 8.2 \\
4.5 & 6.4 & 7.5 \\
4.1 & 5.9 & 8.0 \\
4.3 & 6.2 & 7.9 \\
\bottomrule
\end{tabular}
\end{center}

\vspace{0.5cm}

\textbf{Part A: One-Way ANOVA}

\vspace{0.3cm}

\textbf{R Code:}

\begin{rcode}
# Enter data
program_a <- c(4.2, 3.8, 4.5, 4.1, 4.3)
program_b <- c(6.1, 5.8, 6.4, 5.9, 6.2)
program_c <- c(7.8, 8.2, 7.5, 8.0, 7.9)

# Create data frame
weight_loss <- c(program_a, program_b, program_c)
program <- factor(rep(c("A", "B", "C"), each=5))
data <- data.frame(weight_loss, program)

# Perform ANOVA
anova_model <- aov(weight_loss ~ program, data=data)
summary(anova_model)
\end{rcode}

\vspace{0.3cm}

\textbf{Output:}
\begin{verbatim}
            Df Sum Sq Mean Sq F value   Pr(>F)
program      2 34.200  17.100   268.8 1.08e-10 ***
Residuals   12  0.764   0.064
---
Signif. codes:  0 '***' 0.001 '**' 0.01 '*' 0.05 '.' 0.1 ' ' 1
\end{verbatim}

\vspace{0.3cm}

\textcolor{answergreen}{\textbf{ANSWER: F = 268.85}}

\vspace{0.3cm}

\textbf{Conclusion:} The ANOVA test shows highly significant differences among the three programs ($F(2,12) = 268.85$, $p < 0.001$). We reject the null hypothesis and conclude that at least one program differs in effectiveness.

\vspace{1cm}

\textbf{Part B: Two-Sample t-Test (Program A vs Program C)}

\vspace{0.3cm}

\textbf{R Code:}

\begin{rcode}
# Two-sample t-test (A vs C, equal variances, two-sided)
t_test_result <- t.test(program_a, program_c,
                        var.equal = TRUE,
                        alternative = "two.sided")
print(t_test_result)
\end{rcode}

\vspace{0.3cm}

\textbf{Output:}
\begin{verbatim}
	Two Sample t-test

data:  program_a and program_c
t = -22.601, df = 8, p-value = 2.109e-08
alternative hypothesis: true difference in means is not equal to 0
95 percent confidence interval:
 -4.090863 -3.309137
sample estimates:
mean of x mean of y
     4.18      7.88
\end{verbatim}

\vspace{0.3cm}

\textcolor{answergreen}{\textbf{ANSWER: p-value = 0.0000}}

\vspace{0.3cm}

\textbf{Note:} The exact p-value is $2.109 \times 10^{-8}$, which rounds to 0.0000 when expressed to 4 decimal places.

\vspace{0.3cm}

\textbf{Conclusion:} Program C produces significantly greater weight loss than Program A ($t = -22.60$, $df = 8$, $p < 0.001$). The mean difference is 3.7 kg (95\% CI: 3.31 to 4.09 kg).

\vspace{1cm}

\section*{Summary of Answers}

\begin{center}
\begin{tcolorbox}[colback=answergreen!10,colframe=answergreen,boxrule=1pt,arc=3pt,left=6pt,right=6pt,top=8pt,bottom=8pt]
\textbf{VERIFIED ANSWERS:}\\[0.3cm]
Question 1: \textbf{34 subjects}\\
Question 2A: \textbf{F = 268.85}\\
Question 2B: \textbf{p-value = 0.0000}
\end{tcolorbox}
\end{center}

\vspace{0.5cm}

\noindent\textit{All numerical answers verified using R. See \texttt{verify\_assignment\_answers.R} for complete verification code.}

\end{document}
